\documentclass[11pt]{article}
\usepackage[margin=0.55in]{geometry}
\usepackage{setspace}
\usepackage{booktabs}
\usepackage{longtable}
\usepackage{array}
\usepackage{amsmath, amssymb}
\usepackage{mathtools}
\usepackage{hyperref}
\usepackage{enumitem}

\hypersetup{
  colorlinks=true,
  linkcolor=blue,
  citecolor=blue,
  urlcolor=blue
}

\pdfstringdefDisableCommands{%
  \def\beta{beta}%
  \def\tau{tau}%
  \def\mathbb#1{#1}%
  \def\mathrm#1{#1}%
}

\newcommand{\E}{\mathbb{E}}
\newcommand{\Var}{\mathrm{Var}}
\newcommand{\Cov}{\mathrm{Cov}}
\newcommand{\R}{\mathbb{R}}

\begin{document}
\begin{center}
{\LARGE \textbf{Media Bias ISM}}\\
\vspace{0.15em}
{\normalsize RQ Methods Report}\\
\end{center}
\vspace{0.2em}
\hrule
\vspace{0.5em}
\setstretch{1.05}
\small

\section*{RQ1}
\textbf{Question.} How has media reporting bias toward S\&P 500 firms changed after the 2020 recession shock?

\subsection*{Variable Construction}
For article $a$ of firm $i$ on date $t$, with MTSC probabilities $(p_a^+, p_a^-, p_a^0)$, keep high-confidence observations only:
\[
\max\{p_a^+, p_a^-, p_a^0\} \ge \tau,\quad \tau=0.90\;\text{(default)}.
\]
Define article stance component and firm-day outcome:
\[
s_a = p_a^+ - p_a^-,
\qquad
\text{bias}_{it} = \frac{1}{|\mathcal{A}_{it}|}\sum_{a\in\mathcal{A}_{it}} s_a,
\]
where $\mathcal{A}_{it}$ is the set of qualifying articles for firm $i$ on day $t$.

Controls:
\begin{itemize}[leftmargin=1.5em]
\item $\text{vol}_{it}=|\mathcal{A}_{it}|$ (firm-day article volume)
\item $\text{vix}_t$ (daily market uncertainty)
\end{itemize}

Shock indicator:
\[
\text{Post}_t = \mathbb{1}[t\ge\text{2020-01-01}].
\]

\subsection*{Baseline Model (Firm FE)}
\[
\text{bias}_{it} = \alpha_i + \beta\,\text{Post}_t + \gamma_1\,\text{vol}_{it} + \gamma_2\,\text{vix}_t + \varepsilon_{it}.
\]
$\beta$ is the main estimand: average within-firm post-2020 shift in stance.

\subsection*{Estimator and Inference}
Within transform by firm means, then OLS:
\[
\hat{\theta}=(\tilde{X}^{\top}\tilde{X})^{-1}\tilde{X}^{\top}\tilde{y},\quad
\theta=(\beta,\gamma_1,\gamma_2)^{\top}.
\]
Robust HC1 covariance:
\[
\widehat{\Var}(\hat{\theta}) = \frac{n}{n-K}(\tilde{X}^{\top}\tilde{X})^{-1}
\left(\sum_{j=1}^{n}\hat{u}_j^2\tilde{x}_j\tilde{x}_j^{\top}\right)
(\tilde{X}^{\top}\tilde{X})^{-1}.
\]
Two-sided normal-approx p-values:
\[
p_k = 2\,\Phi(-|t_k|),\quad t_k=\frac{\hat{\theta}_k}{\widehat{\mathrm{SE}}(\hat{\theta}_k)}.
\]

\subsection*{Identification Conditions}
\begin{enumerate}[leftmargin=1.5em]
\item Conditional within-firm parallel trends after controls.
\item Controls are not post-treatment bad controls in the identification sense.
\item \textbf{Stable measurement process:} NewsMTSC score mapping is comparable over time.
\end{enumerate}

\subsection*{Pre-Period Diagnostic and Fallback Trigger}
Compute firm pre means:
\[
\mu_i^{\mathrm{pre}} = \frac{1}{T_i^{\mathrm{pre}}}\sum_{t \in \mathrm{pre}}\text{bias}_{it}.
\]
Track min, max, range, and standard deviation across firms. If range exceeds tolerance (default 0.02), run SCDiD fallback.

\subsection*{SCDiD Fallback (Compact)}
\begin{itemize}[leftmargin=1.5em]
\item For each firm $i$, estimate donor weights from pre period:
\[
\hat{w}_i=\arg\min_w\|y_{\mathrm{pre}}-X_{\mathrm{pre}}w\|_2^2,
\]
then truncate negatives and renormalize.
\item Build post synthetic path $\hat{y}^{\mathrm{syn}}_{it}=X_{\mathrm{post},t}\tilde{w}_i$.
\item Firm effect:
\[
\hat{\tau}_i=\frac{1}{T_i^{\mathrm{post}}}\sum_{t\in\mathrm{post}}(y_{it}-\hat{y}^{\mathrm{syn}}_{it}).
\]
\item Aggregate fallback effect:
\[
\hat{\tau}_{\mathrm{SCDiD}}=\frac{1}{N^*}\sum_{i=1}^{N^*}\hat{\tau}_i,
\quad
\widehat{\mathrm{SE}}=\frac{s_{\hat{\tau}}}{\sqrt{N^*}}.
\]
\end{itemize}

\subsection*{Fallback Eligibility}
A firm is included in fallback only if it has:
\begin{itemize}[leftmargin=1.5em]
\item at least 180 observed pre-period days,
\item at least 180 observed post-period days,
\item sufficient donor coverage for synthetic fitting.
\end{itemize}

\subsection*{Results}
\textbf{Sample and diagnostics.} The estimation panel contains 31,575 firm-day observations across 22 firms (2015-01-01 to 2025-11-12). Pre-period firm means are highly dispersed (range $=0.733$, SD $=0.157$), so the SCDiD fallback is activated.

\begin{table}[ht]
\centering
\caption*{Table RQ1-A: Entity FE Regression (HC1 Robust SE)}
\begin{tabular}{lrrrrrr}
\toprule
Variable & Coef. & SE & t-stat & p-value & CI Low & CI High \\
\midrule
Post & 0.0258 & 0.0052 & 4.9425 & 0.000001 & 0.0155 & 0.0360 \\
Article Volume & -0.0028 & 0.0003 & -9.4603 & 0.000000 & -0.0033 & -0.0022 \\
VIX & 0.0005 & 0.0004 & 1.4886 & 0.1366 & -0.0002 & 0.0012 \\
\bottomrule
\end{tabular}
\end{table}

\begin{table}[ht]
\centering
\caption*{Table RQ1-B: SCDiD Fallback Summary}
\begin{tabular}{lrrr}
\toprule
Metric & Estimate & CI Low & CI High \\
\midrule
Avg. Post Effect & -0.1221 & -0.2034 & -0.0407 \\
SE Across Firms & 0.0415 & -- & -- \\
Firms Used & 16 & -- & -- \\
\bottomrule
\end{tabular}
\end{table}

\textbf{Interpretation.} The baseline FE model indicates a positive post-2020 shift in daily stance. However, the fallback SCDiD estimate is negative, indicating direction sensitivity under strong pre-period imbalance. Report RQ1 as evidence of a structural shift with specification-dependent sign.

\section*{RQ2}
\textbf{Question.} Did the 2020 recession shock affect stock prices differently pre versus post, at the market level?

\subsection*{Design Rationale}
RQ2 is paired with RQ1. RQ1 tests whether media stance shifts after the shock, while RQ2 tests whether firm-level returns also shift. Both are required before claiming a structural media--market linkage in RQ3.

\subsection*{Variable Construction}
For firm $i$ and date $t$, daily return is computed from close prices:
\[
\text{return}_{it}=\frac{P_{it}-P_{i,t-1}}{P_{i,t-1}}.
\]
Shock indicator:
\[
\text{Post}_t=\mathbb{1}[t\ge\text{2020-01-01}].
\]

Controls:
\begin{itemize}[leftmargin=1.5em]
\item $\text{sp500\_return}_t$: absorbs market-wide bull-run and broad index drift.
\item $\text{vix}_t$: absorbs volatility-regime shifts common to all firms.
\end{itemize}

\subsection*{Baseline Model (Firm FE)}
\[
\text{return}_{it}=\alpha_i+\beta\,\text{Post}_t+\gamma_1\,\text{sp500\_return}_t+\gamma_2\,\text{vix}_t+\varepsilon_{it}.
\]
As in RQ1, estimation uses entity fixed effects and HC1 robust standard errors.

\subsection*{Identification and Fallback Rule}
Pre-period firm means are evaluated as in RQ1. If cross-firm pre means are sufficiently equal, FE remains the primary estimate; otherwise SCDiD fallback is triggered.

\subsection*{Results}
\textbf{Sample and diagnostics.} The estimation panel contains 62,667 firm-day observations across 26 firms (2015-12-01 to 2025-11-24). Pre-period means are tightly clustered (range $=0.00193$, SD $=0.00042$), so fallback is \textbf{not} triggered.

\begin{table}[ht]
\centering
\caption*{Table RQ2-A: Entity FE Regression (HC1 Robust SE)}
\begin{tabular}{lrrrrrr}
\toprule
Variable & Coef. & SE & t-stat & p-value & CI Low & CI High \\
\midrule
Post & -0.000344 & 0.000172 & -2.0038 & 0.0451 & -0.000681 & -0.000008 \\
S\&P 500 Return & 1.127576 & 0.010982 & 102.6767 & 0.0000 & 1.106052 & 1.149100 \\
VIX & -0.000001 & 0.000020 & -0.0433 & 0.9654 & -0.000039 & 0.000038 \\
\bottomrule
\end{tabular}
\end{table}

\textbf{Interpretation.} Controlling for index return and VIX, the post-period shift in firm daily returns is negative and marginally significant (about $-0.034\%$ per day). The market-return control is strongly positive and highly significant, while VIX is not statistically significant in this specification.

\end{document}
